%% kapa: exam=1 round=26 year=2015 period=2 form=A subject=accounting subject_ko=회계학 seq=34 number=74
%% source: Q-Net 감정평가사 기출문제 (cst00309) artlSeq=5138001 file=extracted/1008/5138001/제26회_감정평가사 1차_기출/제26회_감정평가사 1차_기출/2교시_감평_A형.pdf
%% conversion: kordoc 4.24.51 -> tools/gen1.py (auto); review=proofread-sonnet; flags=none
\begin{question}{74}
선박을 제조하여 판매하는 (주)감평은 20×5년 초에 영업을 개시하였으며, 제조와 관련된 원가 및 활동에 관한 자료는 다음과 같다.

{\small\centering\begin{tabular}{|c|r|r|r|}\hline
 & \multicolumn{1}{c|}{화물선} & \multicolumn{1}{c|}{유람선} & \multicolumn{1}{c|}{여객선} \\ \hline
\makecell{직접\\재료원가} & ￦60,000 & ￦140,000 & ￦200,000 \\ \hline
\makecell{직접\\노무원가} & 240,000 & 460,000 & 500,000 \\ \hline
\makecell{실제직접\\작업시간} & 1,500시간 & 1,500시간 & 2,000시간 \\ \hline
완성도 & \multicolumn{1}{c|}{60\%} & \multicolumn{1}{c|}{100\%} & \multicolumn{1}{c|}{100\%} \\ \hline
\end{tabular}\par}
\medskip
(주)감평은 직접작업시간을 제조간접원가 배부기준으로 사용하는 정상원가계산제도를 채택하고 있다. 20×5년 제조간접원가예산은 ￦480,000이고 예정 직접작업시간은 6,000시간이다. 20×5년에 발생한 실제 제조간접원가는 ￦500,000이고, 완성된 제품 중 여객선은 고객에게 인도되었다. 제조간접원가 배부차이를 총원가(총원가비례배분법)를 기준으로 조정할 경우 제품원가는?
\begin{choices}
  \item ￦450,000
  \item ￦750,000
  \item ￦756,000
  \item ￦903,000
  \item ￦1,659,000
\end{choices}
\end{question}
\answer{3}
